\documentclass[11pt,letterpaper]{article}
\usepackage[T1]{fontenc}
\usepackage[utf8]{inputenc}
\usepackage{lmodern}
\usepackage[margin=1in,headheight=15pt]{geometry}
\usepackage{amsmath,amssymb,amsthm,mathtools,mathrsfs}
\usepackage{microtype,booktabs,longtable,array,enumitem}
\usepackage[dvipsnames]{xcolor}
\usepackage{fancyhdr,aliascnt,xurl}
\usepackage[colorlinks=true,linkcolor=MidnightBlue,citecolor=MidnightBlue,urlcolor=MidnightBlue]{hyperref}
\usepackage[nameinlink,noabbrev]{cleveref}
\hypersetup{pdftitle={Untransformed Leray Cycles for Radius-Free Hahn Analytic Residues},pdfauthor={Research continuation prepared for Vladimir Reshetnikov},pdfsubject={Surcomplex analysis; Hahn summability; Grothendieck residues; ramified Leray cycles}}
\pagestyle{fancy}\fancyhf{}
\fancyhead[L]{\small Untransformed Leray cycles}
\fancyhead[R]{\small Radius-free Hahn residues}
\fancyfoot[C]{\thepage}
\setlength{\parskip}{0.3em}\setlength{\parindent}{1.3em}
\setlist{nosep,leftmargin=2em}
\allowdisplaybreaks[2]
\newtheorem{theorem}{Theorem}[section]
\newaliascnt{lemma}{theorem}\newtheorem{lemma}[lemma]{Lemma}\aliascntresetthe{lemma}
\newaliascnt{proposition}{theorem}\newtheorem{proposition}[proposition]{Proposition}\aliascntresetthe{proposition}
\newaliascnt{corollary}{theorem}\newtheorem{corollary}[corollary]{Corollary}\aliascntresetthe{corollary}
\theoremstyle{definition}
\newaliascnt{definition}{theorem}\newtheorem{definition}[definition]{Definition}\aliascntresetthe{definition}
\newaliascnt{example}{theorem}\newtheorem{example}[example]{Example}\aliascntresetthe{example}
\theoremstyle{remark}
\newaliascnt{remark}{theorem}\newtheorem{remark}[remark]{Remark}\aliascntresetthe{remark}
\newcommand{\C}{\mathbb C}\newcommand{\R}{\mathbb R}
\newcommand{\N}{\mathbb N}\newcommand{\Z}{\mathbb Z}\newcommand{\Q}{\mathbb Q}
\newcommand{\No}{\mathbf{No}}\newcommand{\SC}{\mathbf{SC}}
\newcommand{\OO}{\mathcal O}\newcommand{\Agerm}{\mathscr A}
\newcommand{\HH}{\mathscr H}\newcommand{\FF}{\mathscr F}
\newcommand{\m}{\mathfrak m}\newcommand{\Res}{\operatorname{Res}}
\newcommand{\supp}{\operatorname{supp}}\newcommand{\CT}{\operatorname{CT}}
\newcommand{\id}{\operatorname{id}}\newcommand{\Jac}{\operatorname{Jac}}
\newcommand{\HInt}{\mathop{\int^{\!H}}\nolimits}
\newcommand{\dd}{\,d}\newcommand{\dz}{\,dz_1\wedge\cdots\wedge dz_n}
\newcommand{\T}{\mathbb T}\newcommand{\D}{\mathbb D}
\newcommand{\wt}{\widetilde}\newcommand{\val}{v}
\newcommand{\smallO}{\mathcal O}
\newcommand{\powers}[1]{\langle #1\rangle_{\N}}
\newcolumntype{L}[1]{>{\raggedright\arraybackslash}p{#1}}
\numberwithin{equation}{section}
\setcounter{tocdepth}{2}

\begin{document}
\hypersetup{pageanchor=false}
\begin{titlepage}
\centering
\vspace*{0.6cm}
{\Large\scshape A research continuation in surcomplex analysis\par}
\vspace{0.85cm}
{\Huge\bfseries Untransformed Leray Cycles\par}
\vspace{0.35cm}
{\LARGE for Radius-Free Hahn Analytic Residues\par}
\vspace{0.7cm}
{\large Ramified finite covers, universal protecting cycles,\par
exact equation-adapted lifts, and support-controlled Stokes\par}
\vspace{0.65cm}
{\large September 21, 2026\par}
\vspace{0.6cm}
\begin{minipage}{0.94\textwidth}
\begin{abstract}
Let $f=(f_1,\ldots,f_n)$ be a convergent complex analytic map germ with an
isolated zero, and let $F=f+E$ be a positive-support Hahn perturbation. The
coefficients of $E$ are individually convergent germs, with no common radius.
The existing surcomplex contour report asks for one admissible cycle realizing
the perturbation residue of the \emph{original} form
$H\,dz_1\wedge\cdots\wedge dz_n/(F_1\cdots F_n)$, rather than a
coordinate-torus representation after a determinant transformation.

We give such a construction. A target monomial arc avoids the discriminant of
$f$. Killing its finite radial monodromy produces one fixed finite cover of a
polyannulus and a holomorphic family of ordinary Leray cycles, ramified only in
one radial parameter. Substitution of a sufficiently large infinitesimal radius
converts that family into a Hahn cycle. The same cycle works for \emph{every}
perturbation with a prescribed positive valuation lower bound, and for every
radius-free numerator. An explicit well-ordered support certificate justifies
all Taylor substitutions, geometric expansions, regroupings, and integrations.

A second theorem constructs an exact equation-adapted lift satisfying
$F_j(\Phi)=t^{m_j\rho}p_j$. It is connected to the protecting cycle by a
pole-free Hahn homotopy. This yields a direct Jacobian-residue normalization,
a finite-cover trace formula, and a separately proved extension to formal
coefficient germs. A mixed six-sheeted singularity and exact symbolic checks
illustrate the construction. The proofs use neither fine-topological limits nor
an Archimedean assumption on the value group.
\end{abstract}
\end{minipage}
\vfill
\begin{minipage}{0.94\textwidth}\small
\textbf{Status and scope.} This article supplies a proof addressing an explicitly
stated research problem in the supplied repository. It does not claim to settle
a named published conjecture. The particular arbitrary-support, radius-free
contour theorem was not found in the targeted literature search; this is not a
priority certification. Classical finite-map, covering, residue, and support
lemmas are credited. The general arguments have not been formally verified or
refereed. Executable checks concern finite examples only.
\end{minipage}
\end{titlepage}
\pagenumbering{roman}\setcounter{page}{2}
\begingroup\footnotesize\setlength{\parskip}{0pt}\tableofcontents\endgroup
\clearpage\pagenumbering{arabic}\hypersetup{pageanchor=true}

\section{The problem and the result}
\label{sec:intro}

\subsection{The repository question}
The repository \emph{Surreal}, inspected at commit
\texttt{39f2be6667ade51bca2b45daa47e289d69c09764}, distinguishes several
non-equivalent analytic categories. Its contour report explicitly separates a
proved determinant-transformed representation from an open request for a
single-cycle representation of the untransformed integrand
\cite{Repo,RepoContours}. Its subsection ``Research problems sharpened by these
results'' asks for an admissible $n$-cycle for arbitrary isolated ordinary $f$
and radius-free positive-support perturbations, with all summation and radius
issues justified. The discussion surrounding its general representation theorem
makes clear that merely reusing its transformed coordinate torus is not an
answer.

The present article addresses exactly that issue. Here ``one cycle'' means one
compact oriented parameter manifold, possibly a finite disjoint union, with one
Hahn parameterization. It does \emph{not} mean a connected coordinate torus in
the original $z$ coordinates. Finite covers and ordinary smooth parameter
manifolds are already legitimate in the coefficientwise chain formalism. We
make the resulting chain category explicit rather than assuming that it is a
category of fine-continuous maps.

\subsection{The coefficient rings}
Fix a nonzero divisible ordered abelian group $\Gamma$, a set rather than a
proper class, and put
\[
 K=\C((t^\Gamma)),\qquad
 \Agerm_n=\C\{z_1,\ldots,z_n\}((t^\Gamma)).
\]
An element of $\Agerm_n$ is a well-ordered Hahn family of convergent coefficient
germs. It is not required to lie in $\OO(D)((t^\Gamma))$ for any common disk
$D$. The latter ring will appear only \emph{after} a justified substitution,
with a fixed ordinary parameter domain. We also use
\[
 \FF_n=\C[[z_1,\ldots,z_n]]((t^\Gamma))
\]
in a separately stated extension theorem. No analytic assertion is silently
transported between these rings.

When $\Gamma$ is a divisible set-sized subgroup of the surreal additive group,
the convention $t^\gamma=\omega^{-\gamma}$ realizes $K$ as a surcomplex
workspace by Conway normal forms. The proofs below are Hahn-field proofs;
they do not use a proper-class topology. Such workspaces are the setting of
the repository's analytic and foundational reports \cite{RepoAnalysis,Repo}.
For $\Gamma=\{0\}$ there is no nonzero positive perturbation, and the relevant
conclusion is the ordinary residue theorem.

Write
\begin{equation}\label{eq:data}
 f\in\C\{z\}^n,\quad f(0)=0,\quad
 \mu=\dim_\C\C\{z\}/(f_1,\ldots,f_n)<\infty,
 \qquad F=f+E.
\end{equation}
For a Hahn family of functions, valuation refers to the least exponent of a
nonzero coefficient function, not to a pointwise minimum over its parameter
domain. Assume the union $S_E$ of the supports of the finitely many $E_j$ is
contained in $[\delta,\infty)$ for a specified $\delta>0$. It is well ordered.
There is no norm bound on its complex coefficients.

\subsection{The independent residue functional}
For $k=(k_1,\ldots,k_n)\in\N^n$, let
\begin{equation}\label{eq:Rk}
 R_{f,k}(h)=\Res_0
 \frac{h(z)\,dz_1\wedge\cdots\wedge dz_n}
 {f_1(z)^{k_1+1}\cdots f_n(z)^{k_n+1}}
\end{equation}
be the ordinary local Grothendieck residue. The ordered list of denominators
fixes the orientation. Equivalently, integrate over a sufficiently small
regular Leray cycle for $f$, oriented by
$d\arg f_1\wedge\cdots\wedge d\arg f_n$.
Classical residue theory supplies this definition and its independence of the
sufficiently small admissible radii \cite{GH,TN}.
Extend $R_{f,k}$ coefficientwise to Hahn families of germs and define
\begin{equation}\label{eq:residue}
 \mathcal R_F(H)
 =\sum_{k\in\N^n}(-1)^{|k|}
 R_{f,k}\bigl(H E_1^{k_1}\cdots E_n^{k_n}\bigr).
\end{equation}
This is the perturbation-residue convention of the repository, not a residue
newly defined to make the desired contour identity tautological. The series is
strongly summable: if $S_H=\supp H$, its support is contained in
$S_H+\powers{S_E}$, and each exponent receives finitely many contributions.
The needed support lemma is recalled in \cref{sec:support}.

\begin{theorem}[Universal untransformed contour]\label{thm:main}
For $f$ and $\delta$ as in \eqref{eq:data}, there exist an ordinary compact
oriented $n$-manifold $T$, a finite covering
$p:T\to\T^n$ of total degree $\mu$, positive integers $m_1,\ldots,m_n$, and
a Hahn parameterization $Z_0:T\to K^n$, with all coordinate supports strictly
positive, such that the following holds.

For every $E\in\Agerm_n^n$ whose support is contained in $[\delta,\infty)$,
and every $H\in\Agerm_n$, the original form has a well-defined coefficientwise
pullback and
\begin{equation}\label{eq:main}
 \boxed{\quad
 \mathcal R_F(H)=\frac{1}{(2\pi i)^n}
 \HInt_{Z_0[T]}
 \frac{H(z)\,dz_1\wedge\cdots\wedge dz_n}{F_1(z)\cdots F_n(z)}.
 \quad}
\end{equation}
The cycle $Z_0[T]$ is chosen from $f$ and $\delta$ alone, independently of
$E$ and $H$. No determinant transformation of $F$ is made.
All sums and all resulting coefficient functions admit one common ordinary
parameter neighborhood of $T$ and an explicit well-ordered Hahn support bound.
\end{theorem}

\begin{theorem}[Equation-adapted lift]\label{thm:adapted-main}
The protecting scale in \cref{thm:main} can be chosen so that, for each $E$,
there is a parameterization $\Phi=Z_0+U$ on the same finite cover with
\begin{equation}\label{eq:adapted}
 F_j(\Phi(q))=t^{m_j\rho}p_j(q),\qquad q\in T.
\end{equation}
Here $\rho>0$, $U$ is strictly smaller than a specified power of $t^\rho$,
and $\Phi$ is uniquely determined within that specified neighborhood.
The straight Hahn homotopy $Z_0+\lambda U$, $0\leq\lambda\leq1$, is
admissible and pole-free for the original form. Thus \eqref{eq:main} also
holds with $\Phi[T]$ in place of $Z_0[T]$.
\end{theorem}

The theorem gives an equation-adapted contour in the precise parameterized
sense \eqref{eq:adapted}. It does not identify its image with the \emph{entire}
class of all surcomplex solutions of $|F_j|=t^{m_j\rho}$.
No such identification is needed for the integral.

\subsection{Where the nontrivial work occurs}
The ordinary map $f$ may have a singular zero and may have a discriminant that
meets many natural target tori. Moreover, replacing a radius-free family by a
common-radius family would be false. Our construction first chooses a monomial
direction in the \emph{target} avoiding the discriminant, and then retains all
sheets over its torus. A finite ramification in one radial variable yields a
fixed parameter manifold. Only afterward is that radial variable replaced by a
Hahn monomial. This order of operations is essential.

The ordinary ramification step belongs to the classical circle of ideas behind
the analytic Abhyankar--Jung theorem; a closely related covering-and-removable-singularity argument appears in Parusi\'nski--Rond \cite[Proposition 2.1]{PR}.
We prove the particular version used here. The proposed contribution is its
combination with a universal protecting scale, arbitrary well-ordered Hahn
supports, radius-free numerator and perturbation families, an exact
support-controlled lift, and equality with the independently defined residue
\eqref{eq:residue}. The classical ingredients themselves are not claimed new.

\section{Support calculus and radius-free substitution}
\label{sec:support}

\subsection{Strong summability}
\begin{definition}
A family $(a_i)_{i\in I}$ of Hahn series, possibly with coefficients in a
complex vector space, is \emph{strongly summable} if the union of its supports
is well ordered and, at each exponent, only finitely many indexed terms are
nonzero. Its sum is defined coefficientwise. These conditions, not convergence
of partial sums, are used throughout.
\end{definition}

For a subset $P$ of the positive cone, write
\[
 \powers P=\{p_1+\cdots+p_r:r\geq0,\ p_j\in P\},
\]
including $0$ from the empty sum. The following classical fact is often called
the Neumann support lemma \cite{Neumann}.

\begin{lemma}[Positive support lemma]\label{lem:neumann}
If $P\subset\Gamma_{>0}$ is well ordered, then $\powers P$ is well ordered.
Each $\gamma\in\Gamma$ is the sum of only finitely many finite ordered words
with letters in $P$. If $S$ is any well-ordered subset of $\Gamma$, then
$S+\powers P$ is well ordered, and a fixed exponent has only finitely many
representations as an element of $S$ plus such a word.
\end{lemma}
\begin{proof}
For completeness, one route to the nontrivial assertion is Higman's classical
lemma for words on a well-quasi-ordered alphabet \cite{Higman}. A well-order is
a well-quasi-order. If one word embeds into another with every matched letter
less than or equal to its partner, its sum is less than or equal to the other
sum. Equality forces identical words: every omitted letter is strictly
positive, and every strict increase is positive. A descending sequence of word
sums, or infinitely many distinct words with one fixed sum, therefore
contradicts Higman's lemma. This proves the first two assertions.

For two well-ordered subsets $A,B$ of an ordered group, $A+B$ is well ordered:
from any alleged descending sequence, extract successively nondecreasing
subsequences of its $A$ and $B$ entries. A fixed sum has only finitely many
representations, because infinitely many distinct $A$ entries would have an
increasing subsequence and hence force a decreasing sequence of $B$ entries.
Apply this with $A=S$ and $B=\powers P$, and use the finite-word assertion.
\end{proof}

\begin{remark}\label{rem:notcofinal}
The lemma does \emph{not} say that only finitely many elements of
$\powers P$ lie below a given bound. That stronger statement is generally
false, even for $P=\{\rho\}$. In a higher-rank group, all $r\rho$ can lie
below a fixed positive $\gamma$. The proofs use finiteness at an \emph{exact}
exponent, never finiteness of an initial segment.
\end{remark}

\subsection{Coefficients on a fixed parameter manifold}
Let $X$ be an ordinary complex manifold. Write
$\HH(X)=\OO(X)((t^\Gamma))$ for families whose coefficient functions are
holomorphic on this one fixed $X$, and write $\HH(X)_+$ for those with
strictly positive support. We use this notation only for the displayed domain;
no assertion about an unrestricted sheaf on arbitrary disconnected covers is
needed. For a compact real manifold $T\subset X$, coefficientwise integration
of top-degree smooth forms is the linear map
\begin{equation}\label{eq:Hint}
 \HInt_T\sum_{\gamma}t^\gamma\alpha_\gamma
 =\sum_\gamma t^\gamma\int_T\alpha_\gamma.
\end{equation}
It cannot enlarge the support and commutes with strongly summable families.

\begin{lemma}[Radius-free substitution]\label{lem:substitution}
Suppose $W=(W_1,\ldots,W_n)\in\HH(X)_+^n$ and the union $P_W$ of its
coordinate supports is well ordered. For
$H=\sum_{\gamma\in S_H}t^\gamma h_\gamma\in\FF_n$, the expression
\begin{equation}\label{eq:substitution}
 H(W)=\sum_{\gamma\in S_H}\sum_{\alpha\in\N^n}
 t^\gamma [z^\alpha]h_\gamma\,W^\alpha
\end{equation}
is strongly summable and belongs to $\HH(X)$, with support in
$S_H+\powers{P_W}$. It respects sums, products, and differentiation with
respect to the ordinary parameters on $X$. It applies in particular to
$H\in\Agerm_n$, without a common convergence radius.
\end{lemma}
\begin{proof}
Each fully expanded monomial selects finitely many positive letters from
$P_W$, together with one exponent from $S_H$. At a fixed total exponent,
\cref{lem:neumann} leaves finitely many words. There are finitely many ways to
assign their letters to the $n$ coordinates and to the finitely many factors.
Thus the coefficient is a finite sum of holomorphic functions on $X$.
The same argument justifies all regroupings proving the ring identities.
Ordinary differentiation of a coefficient function leaves its exponent
unchanged. Termwise differentiation follows because at each exponent all
algebraic sums are finite. No convergence of the formal $h_\gamma$ was used.
\end{proof}

\begin{corollary}[Radial evaluation]\label{cor:radial}
Let $W(q,s)=\sum_{\ell\geq1}s^\ell W_\ell(q)$ be holomorphic on
$X\times\D_\epsilon$, with $W(q,0)=0$. For any $\rho>0$, the substitution
$s=t^\rho$ gives an element of $\HH(X)_+^n$. Every radius-free $H(W)$ is
legitimate. Its coefficients are formed from finite Taylor jets, even when
there is no ordinary $\epsilon$ that works simultaneously for all
$h_\gamma(W(q,s))$.
\end{corollary}

\subsection{A nonlinear support lemma}
The following formal implicit-function statement will supply the equation-adapted cycle. The coefficient ring need not be a field.

\begin{lemma}[Positive implicit equation]\label{lem:implicit}
Let $R$ be a commutative complex algebra and let
$B(V)=(B_1(V),\ldots,B_n(V))$ be a formal power series in finitely many
variables $V$, with Hahn-series coefficients in $R((t^\Gamma))$.
Suppose the union of the supports of \emph{all} coefficients of $B$ is a
well-ordered subset $W\subset\Gamma_{>0}$. Then
\[
 V=B(V)
\]
has exactly one solution with positive coordinate supports. Its support is
contained in $\powers W\setminus\{0\}$, and it is obtained by strongly
summable finite rooted-tree expansions. Each output coefficient uses only
finitely many coefficient operations in $R$.
\end{lemma}
\begin{proof}
Expand repeated substitution into rooted ordered trees. A vertex selects a
coefficient of some component of $B$; the multi-index of its monomial specifies
the components and number of its children. Nullary vertices select constant
terms. Assign to each vertex a positive exponent from the selected coefficient.
The weight of a tree is the sum of its vertex exponents.

At a fixed weight, \cref{lem:neumann} permits only finitely many ordered words
of vertex exponents. In particular the number of vertices is bounded for that
weight. There are finitely many ordered rooted trees with that many vertices,
finitely many component colorings, and finitely many compatible multi-indices:
a vertex cannot have more children than the total number of other vertices.
Therefore the tree family is strongly summable. Joining trees below one root
proves $V=B(V)$ by permitted regrouping.

For uniqueness, if $V\ne V'$ are positive solutions, take the smallest
exponent at which a coordinate differs. In the difference
$B(V)-B(V')$, every term contains such a difference multiplied by a positive
coefficient exponent and zero or more positive coordinate exponents. None can
contribute at that smallest exponent, a contradiction.
\end{proof}

This proof does not identify a support-defined sum with a limit of Picard
iterates. In particular it remains valid when the multiples of the least
coefficient valuation are not cofinal in $\Gamma$.

\section{A fixed ramified family of ordinary Leray cycles}
\label{sec:ordinary}

This section is ordinary complex analysis. Its radial variable $s$ is an
ordinary complex number until \cref{sec:universal}. The distinction prevents
an unjustified interchange of analytic convergence and Hahn summability.

\subsection{Finite representatives and discriminants}
We use the classical finite mapping theorem in the following standard form.
If $f:(\C^n,0)\to(\C^n,0)$ has an isolated zero, one can choose bounded
representatives $U,V$ of the source and target, with $f:U\to V$ proper and
finite, $f^{-1}(0)=\{0\}$, and degree
$\mu=\dim_\C\C\{z\}/(f)$. There is a nonzero analytic germ
$\Delta(w)$ such that outside $\{\Delta=0\}$ this map is an unramified
covering of degree $\mu$. These are ordinary finite-map facts, not consequences
asserted for a Hahn analytic space \cite[Chapters II--III]{GR}.

One may see the geometric origin of the representatives by taking a small
closed source ball whose boundary avoids $f^{-1}(0)$, then a small target
polydisk whose inverse image stays away from that boundary. Properness follows.
The fibers are finite because a positive-dimensional compact analytic subset
of a domain in $\C^n$ cannot occur. The finite mapping theorem supplies the
local degree and analyticity of the branch locus. A local defining function,
or a nonzero function vanishing on that branch locus, serves as $\Delta$.
If $f$ is already unramified near $0$, one may use $\Delta=1$.

\begin{lemma}[A discriminant-avoiding monomial direction]\label{lem:weights}
For a nonzero convergent germ $\Delta(w_1,\ldots,w_n)$ there are positive
integers $q_1,\ldots,q_n$, an integer $d\geq0$, a multi-index $\alpha$, and
$c\ne0$ such that on a fixed polyannulus
$A=\{a<|\xi_j|<b\}$ around $\T^n$,
\begin{equation}\label{eq:delta}
 \Delta(\tau^{q_1}\xi_1,\ldots,\tau^{q_n}\xi_n)
 =\tau^d\bigl(c\xi^\alpha+\tau R(\xi,\tau)\bigr).
\end{equation}
After shrinking the $\tau$ disk and, if necessary, the polyannulus, the factor
in parentheses is nowhere zero on a neighborhood of
$\overline A\times\{0\}$. In particular the displayed target points avoid
the branch locus for every sufficiently small $\tau\ne0$.
\end{lemma}
\begin{proof}
Let $d_0$ be the least total degree of a nonzero term of $\Delta$. Choose
positive rational weights very close to $1$, generic on the finitely many
monomials of total degree $d_0$. Exactly one of those monomials has least
weight. The weights can be chosen so close to $1$ that every monomial of total
degree at least $d_0+1$ has larger weight than all relevant degree-$d_0$
monomials. Clearing denominators gives positive integer weights and a unique
least-weight monomial $c w^\alpha$.

Substitution in the convergent series gives \eqref{eq:delta}; all other
weighted degrees exceed $d$ by a positive integer. The remainder is
holomorphic for bounded $\xi$ and small $\tau$. On a closed polyannulus,
$c\xi^\alpha$ is bounded away from zero. Compactness therefore gives uniform
nonvanishing for small $\tau$.
\end{proof}

This argument uses a \emph{target} monomial direction. It does not require the
lowest weighted homogeneous parts of the original $f_j$ to form an isolated
complete intersection in the source.

\subsection{Removing radial monodromy}
\begin{theorem}[Ramified Leray family]\label{thm:ordinary}
For $f$ as in \eqref{eq:data}, there exist a connected polyannulus $A$ around
$\T^n$, a finite unramified holomorphic cover $p:X\to A$ of total degree
$\mu$, positive integers $m_j$, and a holomorphic map
\[
 Z:X\times\D_\epsilon\longrightarrow\C^n
\]
satisfying
\begin{equation}\label{eq:ordinaryZ}
 Z(q,0)=0,\qquad f_j(Z(q,s))=s^{m_j}p_j(q).
\end{equation}
For $s\ne0$, this map identifies the fiber of $p$ over $\xi$ with all
$\mu$ inverse images in $U$ of $(s^{m_j}\xi_j)_j$. The restriction
$T=p^{-1}(\T^n)$ is a compact oriented real $n$-manifold without boundary,
and $Z(T,s)$ is the ordinary Leray cycle for target radii $|s|^{m_j}$.

One can take $m_j=Nq_j$, where $q_j$ are as in \cref{lem:weights} and $N$
divides $\operatorname{lcm}(1,\ldots,\mu)$.
\end{theorem}
\begin{proof}
Pull the finite map $f:U\to V$ back by
\[
 A\times\D^*\longrightarrow V,\qquad
 (\xi,\tau)\longmapsto(\tau^{q_j}\xi_j)_j.
\]
The discriminant-avoidance lemma makes the pullback an unramified holomorphic
cover of degree $\mu$. Its monodromy is a permutation representation of
\[
 \pi_1(A\times\D^*)\simeq\Z^n\times\Z.
\]
Let $\sigma$ be the permutation induced by the last, radial generator, and
let $N$ be its order. The substitution $\tau=s^N$ kills that monodromy.
The remaining covering is pulled back from a finite cover $p:X\to A$:
its representation is trivial on the last factor. The classification of
covering spaces gives an isomorphism with $X\times\D^*$. This isomorphism is
holomorphic, since locally it is an isomorphism between holomorphic local
inverses over the same base. No global labeling of all sheets over $A$ is
required. The order bound on $N$ is the elementary order bound for a
permutation of $\mu$ letters.

The source-coordinate projection now gives a holomorphic map
$Z:X\times\D^*\to U$ satisfying \eqref{eq:ordinaryZ} for $s\ne0$.
It is bounded, because $U$ was chosen bounded. The removable-singularity
theorem with holomorphic parameters extends it across $s=0$. This can be
proved in local charts by the Cauchy integral formula in $s$; boundedness kills
all negative Laurent coefficients and the remaining coefficients are
holomorphic in $q$. Extensions on overlaps agree. By continuity,
$f(Z(q,0))=0$, whence $Z(q,0)=0$.

For fixed $s\ne0$, every inverse image of the target torus occurs once in
this finite cover. The source-coordinate map is an embedding on $T$: equality
of source points implies equality of target coordinates, and then equality
of the sheet. Orient $T$ so that
$p^*(d\arg\xi_1\wedge\cdots\wedge d\arg\xi_n)$ is positive. Multiplying
each target coordinate by $s^{m_j}$ only adds a constant phase; the induced
orientation is exactly the Leray orientation. Finally, a finite cover of a
compact torus is compact and has finitely many components, each a torus.
\end{proof}

The parameter space $X$ is allowed to be disconnected. This avoids falsely
identifying a potentially nontrivial finite covering with a list of globally
single-valued branches on the original polyannulus.

\subsection{Laurent identities before Hahn substitution}
For $h\in\C\{z\}$ and $k\in\N^n$, define the $n$-form on $X$
\begin{equation}\label{eq:ordinaryform}
 \beta_{h,k}(s)=
 \frac{h(Z(q,s))\,d_qZ_1(q,s)\wedge\cdots\wedge d_qZ_n(q,s)}
 {s^{\sum_jm_j(k_j+1)}p_1(q)^{k_1+1}\cdots p_n(q)^{k_n+1}}.
\end{equation}
Here $d_q$ differentiates only the angular parameters, not $s$.

\begin{lemma}[Coefficientwise constant-period identity]\label{lem:constant}
For each fixed $h,k$, \eqref{eq:ordinaryform} has a convergent Laurent
expansion in $s$ with finitely many negative powers, on a neighborhood of $T$
and for sufficiently small $s$. Its coefficientwise integral satisfies
\begin{equation}\label{eq:constantperiod}
 \frac{1}{(2\pi i)^n}\int_T\beta_{h,k}(s)=R_{f,k}(h).
\end{equation}
Thus every nonconstant Laurent coefficient has integral zero, and the constant
coefficient has integral $(2\pi i)^nR_{f,k}(h)$.
\end{lemma}
\begin{proof}
For each fixed $h$, compactness of $T$ and $Z(q,0)=0$ give a sufficiently small
ordinary $s$ disk in which $h(Z(q,s))$ is holomorphic on one neighborhood of
$T$. The other numerator factors are holomorphic; the denominator contributes
only the displayed finite negative power of $s$, since the $p_j$ are units.
Ordinary compact-uniform convergence justifies integration of this one Laurent
expansion. For positive real $s$ sufficiently small, $Z(T,s)$ is an ordinary
Leray cycle and its integral is the residue. Holomorphicity in the punctured
$s$ disk, or uniqueness of a Laurent expansion, proves the coefficient
identity. The small disk may depend on $h$ and $k$; no simultaneous disk for
an infinite Hahn family is asserted.
\end{proof}

\section{The universal protecting cycle and its residue}
\label{sec:universal}

\subsection{Choice of a protecting radius}
Fix the ordinary data of \cref{thm:ordinary}. Set
\[
 m_* =\max_jm_j,\qquad M=\sum_jm_j.
\]
Choose $\rho\in\Gamma_{>0}$ such that
\begin{equation}\label{eq:protect}
 m_*\rho<\delta.
\end{equation}
Divisibility ensures such a choice, for example
$\rho=\delta/(m_*+1)$. Notice the direction: smaller positive exponent means
a \emph{larger} infinitesimal monomial $t^\rho$. The contour must be large
enough, in valuation terms, to protect the whole perturbation.

Expand the ordinary family
$Z(q,s)=\sum_{\ell\geq1}s^\ell Z_\ell(q)$ and define
\begin{equation}\label{eq:Z0}
 Z_0(q)=\sum_{\ell\geq1}t^{\ell\rho}Z_\ell(q).
\end{equation}
This is a legitimate Hahn parameterization on the fixed ordinary $X$, with
\[
 f_j(Z_0)=t^{m_j\rho}p_j.
\]
All coordinate supports are in $\N_{>0}\rho$.

\subsection{An explicit support certificate}
For the particular perturbation $E$ under consideration, put
\begin{equation}\label{eq:P0}
 P_0=\{\rho\}\ \cup\ \bigcup_{j=1}^n(S_E-m_j\rho),
 \qquad \mathcal M_0=\powers{P_0}.
\end{equation}
Finite unions and translations of well-ordered subsets are well ordered.
Condition \eqref{eq:protect} makes every element of $P_0$ positive.

\begin{proposition}[Admissibility on the protecting cycle]\label{prop:admissible}
For every $H\in\Agerm_n$, the pullback of the original form by $Z_0$ belongs to
the Hahn family of holomorphic $n$-forms on $X$ and has support contained in
\begin{equation}\label{eq:support-original}
 S_H-M\rho+\mathcal M_0.
\end{equation}
The assertion includes the fully expanded family before any Taylor,
perturbation, or radial index is regrouped. Each coefficient is a finite sum
of ordinary holomorphic forms on $X$.
\end{proposition}
\begin{proof}
Radius-free substitution gives
\[
 H(Z_0)\in\HH(X),\qquad
 \supp H(Z_0)\subset S_H+\N\rho,
\]
and similarly $\supp E_j(Z_0)\subset S_E+\N\rho$. Define
\[
 Q_j=\frac{E_j(Z_0)}{t^{m_j\rho}p_j}.
\]
Its support is contained in $S_E-m_j\rho+\N\rho$, a positive subset of
$\mathcal M_0$. Hence
\begin{equation}\label{eq:inverse}
 \frac1{F_j(Z_0)}
 =t^{-m_j\rho}p_j^{-1}\sum_{r\geq0}(-Q_j)^r
\end{equation}
is a strongly summable identity in $\HH(X)$. The coefficients $p_j^{-1}$
are ordinary holomorphic units on the same $X$.

Each angular derivative of $Z_0$ still has support in $\N_{>0}\rho$.
Multiplication with \eqref{eq:inverse} gives the bound
\eqref{eq:support-original}. For the stronger indexed-family assertion, expand
a contribution with $k_j$ factors from $E_j$. Its exponent has the form
\begin{equation}\label{eq:exponent-accounting}
 \gamma-M\rho+
 \sum_{j=1}^n\sum_{a=1}^{k_j}(e_{j,a}-m_j\rho)+r\rho,
 \qquad \gamma\in S_H,\ e_{j,a}\in S_E,\ r\geq0.
\end{equation}
Every letter after the fixed shift belongs to $P_0$. At an exact total
exponent, \cref{lem:neumann} permits only finitely many words and hence finitely
many $k$, perturbation coefficients, and radial degrees. For a fixed radial
degree, Taylor substitution in $Z(q,s)$ uses only finitely many Taylor
monomials and finitely many radial coefficients, since $Z(q,0)=0$.
This proves the required local finiteness before regrouping.
\end{proof}

\subsection{Proof of the universal contour theorem}
\begin{proof}[Proof of \cref{thm:main}]
Choose $T,p,Z_0$ as above. Expand the reciprocal factors in
\eqref{eq:inverse}. The proposed normalized integral is
\begin{equation}\label{eq:expandedintegral}
 \frac1{(2\pi i)^n}
 \sum_{k\in\N^n}(-1)^{|k|}
 \HInt_T
 \frac{(HE_1^{k_1}\cdots E_n^{k_n})(Z_0)
 \,dZ_{0,1}\wedge\cdots\wedge dZ_{0,n}}
 {\prod_j f_j(Z_0)^{k_j+1}}.
\end{equation}
The entire expanded family is strongly summable by
\cref{prop:admissible}, so integration and all regroupings are legal.

Fix one selection of Hahn coefficients from $H$ and the finitely many
perturbation factors in one $k$ term. Their ordinary product is a single
convergent germ $h$. Before replacing $s$ by $t^\rho$, its angular pullback
is exactly \eqref{eq:ordinaryform}. By \cref{lem:constant}, integration kills
all nonconstant radial Laurent coefficients and returns the ordinary residue
$R_{f,k}(h)$. Multiplying by the chosen Hahn monomial and summing the
strongly summable family gives precisely \eqref{eq:residue}.

Only $f$ and $\delta$ entered the construction of $T,p,Z_0$. The support
certificate depends on the support of the particular $E$, but the contour does
not. The numerator was arbitrary throughout. This proves all assertions.
\end{proof}

\begin{corollary}[Cancellation of auxiliary scales]\label{cor:cancel}
Although the pulled-back form may involve exponents in the larger set
$S_H-M\rho+\mathcal M_0$, its integral has support contained in
$S_H+\powers{S_E}$. All auxiliary scale contributions cancel after angular
integration. The value is independent of the chosen discriminant, target
weights, ramification, covering identification, and protecting scale.
\end{corollary}
\begin{proof}
The resulting value equals the independent series \eqref{eq:residue}, whose
support and definition do not involve those choices.
\end{proof}

This proves independence of the \emph{period}. It does not assert that every
pair of differently scaled cycles is joined by a pole-free homotopy in some
unspecified larger chain category.

\subsection{Why the missing common radius causes no gap}
Two different finiteness principles appear, and they must not be confused.
For a \emph{fixed ordinary coefficient product}, ordinary compactness supplies
a small $s$ disk and proves its Laurent period identity. For a \emph{fixed Hahn
exponent}, the positive support lemma supplies a finite list of coefficient
products and radial terms. Thus no step asks for one ordinary disk serving
all coefficient germs. The contour's common parameter domain comes from the
finite map $f$ and its covering, not from an intersection of infinitely many
convergence neighborhoods of the perturbation coefficients.

The conclusion is stronger than a representation separately constructed for
one numerator or one perturbation. It gives a common cycle for an entire
valuation ball of positive perturbations, even when the coefficient germs
throughout that ball have arbitrarily small radii.

\section{An exact equation-adapted Hahn lift}
\label{sec:lift}

\subsection{Uniform ordinary pole and separation bounds}
The following bounds depend only on the finite map and its ramified family.
They do not depend on any Hahn coefficient germ.

\begin{lemma}[Finite radial bounds]\label{lem:bounds}
After shrinking the ordinary polyannulus and the ordinary $s$ disk around the
same compact $T$, there are integers $b,c\geq0$ such that:
\begin{equation}\label{eq:Jbound}
 (D_zf(Z(q,s)))^{-1}=s^{-b}A(q,s)
\end{equation}
with $A$ holomorphic on the fixed product neighborhood; and for any distinct
$q,q'\in T$ with $p(q)=p(q')$, at least one coordinate of
$Z(q,s)-Z(q',s)$ has a nonzero Taylor coefficient of degree at most $c$.
If $\mu=1$, take $c=0$.
\end{lemma}
\begin{proof}
The determinant $a(q,s)=\det D_zf(Z(q,s))$ is holomorphic across $s=0$ and
is nowhere zero for $s\ne0$. Locally near $(q_0,0)$, Weierstrass preparation
in $s$ factors it as a unit times a monic polynomial in $s$. Every root of
that polynomial must be $0$, because no zero with $s\ne0$ is allowed in the
product neighborhood. Thus locally $a=s^r u$ with $u$ a holomorphic unit.
The exponent is locally constant along $X\times\{0\}$. A finite collection
of neighborhoods covering $T$ gives a common finite pole bound for the
inverse matrix, via its adjugate. Shrinking to a common collar of $T$ and an
$s$ disk gives \eqref{eq:Jbound}.

For separation, consider the off-diagonal fiber product
$T\times_{\T^n}T\setminus\operatorname{diag}(T)$. It is compact because
it is a union of components of a finite cover. At each of its points the
two holomorphic source-coordinate vectors differ for $s\ne0$, so at least
one Taylor coefficient of their difference is nonzero. Nonvanishing of that
coefficient persists on an ordinary neighborhood. A finite subcover gives a
uniform upper bound $c$ on a degree at which a difference is detected.
\end{proof}

Choose
\begin{equation}\label{eq:Lrho}
 L=\max\{b,c,m_*\}+1,\qquad
 0<\rho,\qquad (L+b)\rho<\delta.
\end{equation}
For example, $\rho=\delta/[2(L+b+1)]$ is available in the divisible group.
Reconstruct $Z_0$ using this possibly smaller exponent. The protecting-cycle
theorem remains valid. Put
\begin{equation}\label{eq:P}
 P=\{\rho\}\cup\bigl(S_E-(L+b)\rho\bigr),
 \qquad \mathcal M=\powers P.
\end{equation}
This is a well-ordered positive generating set. The earlier generating set
$P_0$ is contained in $\mathcal M$, because $L+b\geq m_j$ for every $j$.

\subsection{The scaled implicit equation}
Write $s=t^\rho$ only as an algebraic abbreviation from now on, and set
\[
 J_0=D_zf(Z_0),\qquad \Phi=Z_0+s^L V.
\]
The equation $F(\Phi)=f(Z_0)$ is equivalent to
\begin{align}
 V&=B(V),\label{eq:Bequation}\\
 B(V)&=-s^{-L}J_0^{-1}\Bigl(
 f(Z_0+s^LV)-f(Z_0)-s^LJ_0V+E(Z_0+s^LV)\Bigr).
 \label{eq:Bdefinition}
\end{align}
This is a formal power series in $V$ with coefficients in $\HH(X)$.

\begin{proposition}[Support-controlled lift]\label{prop:lift}
Equation \eqref{eq:Bequation} has a unique positive solution $V$. The
correction $U=s^LV$ satisfies
\begin{equation}\label{eq:Usupport}
 \supp U\subset L\rho+(\mathcal M\setminus\{0\}),\qquad
 \supp U\subset S_E-b\rho+\mathcal M.
\end{equation}
In particular $v(U)>L\rho$ and $v(U)\geq\delta-b\rho$ whenever $U\ne0$.
The map $\Phi=Z_0+U$ satisfies \eqref{eq:adapted}, on the whole fixed
ordinary parameter domain $X$. No further shrinking depending on $E$ occurs.
\end{proposition}
\begin{proof}
In a coefficient of total $V$ degree $r\geq2$ from the ordinary nonlinear
part of \eqref{eq:Bdefinition}, the possible exponents are contained in
\[
 ((r-1)L-b)\rho+\N\rho.
\]
Since $L>b$, these are positive and belong to $\mathcal M$. In a coefficient
of total degree $r\geq0$ from the perturbation, the support is contained in
\[
 S_E-(L+b)\rho+rL\rho+\N\rho,
\]
again a positive subset of $\mathcal M$. These assertions follow by Taylor
expansion and \eqref{eq:Jbound}; all coefficient functions are holomorphic
on the same $X$ by \cref{lem:substitution}. The union of the supports of all
coefficients of $B$ is therefore a well-ordered subset of
$\mathcal M\setminus\{0\}$. Apply \cref{lem:implicit}.

The first support bound in \eqref{eq:Usupport} follows immediately. For the
second, every finite rooted tree has a nullary vertex. The constant coefficient
of $B$ comes entirely from $-s^{-L}J_0^{-1}E(Z_0)$ and has support in
$S_E-(L+b)\rho+\N\rho$. Remove one such vertex exponent from the weight
of the tree; all remaining exponents belong to $\mathcal M$. Thus
$\supp V\subset S_E-(L+b)\rho+\mathcal M$, giving the second bound.
This also proves that the correction is zero when $E=0$.

All substitutions used to derive \eqref{eq:Bequation} are justified by the
same support bounds. Reversing the algebra therefore gives $F(\Phi)=f(Z_0)$
exactly, not modulo a truncation. Uniqueness in the positive $V$ ball follows
from \cref{lem:implicit}; equivalently, it is uniqueness among corrections
with $v(U)>L\rho$.
\end{proof}

\subsection{Nonsingularity and preservation of the finite cover}
\begin{corollary}[Regular protected sheets]\label{cor:sheets}
The matrix $D_zF(\Phi)$ is invertible in the Hahn matrix algebra on $X$.
For fixed ordinary $\xi\in\T^n$, the $\mu$ values $\Phi(q)$ with
$p(q)=\xi$ are pairwise distinct. In local holomorphic coordinates on $X$,
the angular derivative of $\Phi$ has full rank over the Hahn coefficient ring.
\end{corollary}
\begin{proof}
Write
\[
 J_0^{-1}D_zF(\Phi)=I+J_0^{-1}(D_zF(\Phi)-J_0).
\]
The ordinary part of the difference contains at least one factor $U$, so
its exponents after multiplication by $J_0^{-1}$ are strictly above
$(L-b)\rho>0$. The perturbation derivative contributes exponents at least
$\delta-b\rho>0$. A matrix geometric series gives the inverse of the
right-hand side. This proves invertibility of $D_zF(\Phi)$.

For two distinct sheets, \cref{lem:bounds} detects a nonzero coefficient of
$Z_0(q)-Z_0(q')$ at an exponent at most $c\rho$. The correction difference
has all exponents strictly above $L\rho>c\rho$, so it cannot cancel that
coefficient. Finally, differentiating $F_j(\Phi)=t^{m_j\rho}p_j$ gives
\[
 D_zF(\Phi)\,d\Phi
 =\operatorname{diag}(t^{m_1\rho},\ldots,t^{m_n\rho})\,dp.
\]
The right-hand matrix has invertible determinant in local coordinates, because
$p$ is an unramified cover and the monomials are nonzero units in $K$.
\end{proof}

These are protected-sheet statements. They neither count every solution in an
unspecified larger surreal domain nor assume that the displaced zero scheme is
reduced. Singular zeros are allowed; the contour itself lies over regular
nonzero target values.

\section{Hahn Stokes, periods, and normalization}
\label{sec:stokes}

\subsection{The pole-free homotopy}
Let $\Phi_\lambda=Z_0+\lambda U$ for $0\leq\lambda\leq1$, with $U$ from
\cref{prop:lift}. The parameter $\lambda$ is ordinary real; it does not index
a limit in $K$.

\begin{lemma}[Admissible straight homotopy]\label{lem:homotopy}
For each $j$,
\[
 F_j(\Phi_\lambda)=t^{m_j\rho}p_j(1+Q_{j,\lambda}),
\]
where $Q_{j,\lambda}$ has positive support in $\mathcal M$ and coefficients
smooth on $X\times[0,1]$. All Taylor, inverse, differentiation, and pullback
operations for the original form are strongly summable there. Its total
pullback has support in $S_H-M\rho+\mathcal M$.
\end{lemma}
\begin{proof}
The difference $f_j(Z_0+\lambda U)-f_j(Z_0)$ contains a factor $U$, so its
support is contained in $L\rho+(\mathcal M\setminus\{0\})$. Division by
$t^{m_j\rho}p_j$ leaves positive support because $L>m_j$. The perturbation
part has support in $S_E+\mathcal M$; division leaves positive support since
$S_E-m_j\rho\subset\mathcal M\setminus\{0\}$. Expand $(1+Q)^{-1}$
geometrically. The positive support lemma controls the full indexed family.
At each exact exponent the dependence on $\lambda$ is polynomial, because
only finitely many Taylor and geometric-series terms contribute. Ordinary
parameter differentiation cannot enlarge support. This proves the assertion
for the full pullback, including its $d\lambda$ terms.
\end{proof}

\begin{theorem}[Untransformed homotopy invariance]\label{thm:homotopy}
For every $H\in\Agerm_n$,
\begin{equation}\label{eq:homotopy-integral}
 \HInt_{\Phi[T]}\frac{H\dz}{F_1\cdots F_n}
 =\HInt_{Z_0[T]}\frac{H\dz}{F_1\cdots F_n}.
\end{equation}
Consequently \cref{thm:adapted-main} holds with the support and uniqueness
bounds given in \cref{prop:lift}.
\end{theorem}
\begin{proof}
A meromorphic holomorphic top-degree form in $n$ complex variables is closed
where its denominators are invertible: differentiating its coefficient gives
an extra $dz_j$, which wedges to zero with $dz_1\wedge\cdots\wedge dz_n$.
This is also a formal differential-algebra identity. The support-justified
chain rule carries it to the total pullback under $\Phi_\lambda$.

By \cref{lem:homotopy}, that pullback is a strongly summable family of smooth
closed $n$-forms on the compact ordinary cylinder $T\times[0,1]$. Apply
ordinary Stokes to each coefficient. Since $\partial T=\varnothing$, the only
boundary contributions are the two endpoint cycles, with opposite signs.
Summing their zero differences coefficientwise proves
\eqref{eq:homotopy-integral}. Combine with \cref{thm:main}.
\end{proof}

The ordinary pullback Jacobian $d\Phi_1\wedge\cdots\wedge d\Phi_n$ is,
of course, present. It is the change-of-parameters factor of an integral. It
is not an auxiliary determinant $\det A$ from replacing the equation system
$F$ by $AF$. Neither the integrand nor its denominator list has been replaced.

\subsection{A direct Jacobian-residue conservation law}
Put $J_F=\det D_zF$.
\begin{theorem}[Degree normalization]\label{thm:normalization}
Under the hypotheses of \cref{thm:main},
\begin{equation}\label{eq:jacobian}
 \mathcal R_F(J_F)=\mu.
\end{equation}
More generally, for $G\in\Agerm_n$ in target variables $w$,
\begin{equation}\label{eq:target-normalization}
 \mathcal R_F\bigl(G(F)J_F\bigr)=\mu\,G(0).
\end{equation}
Here $G(F)$ is expanded around the ordinary map $f$, in the positive
perturbation $E$, so it is a well-defined radius-free Hahn germ.
\end{theorem}
\begin{proof}
Use the adapted cycle and the identity
\[
 \Phi^*\left(\frac{J_F\dz}{F_1\cdots F_n}\right)
 =\bigwedge_{j=1}^n\frac{d(F_j\circ\Phi)}{F_j\circ\Phi}
 =p^*\left(\frac{dw_1}{w_1}\wedge\cdots\wedge\frac{dw_n}{w_n}\right).
\]
The normalized integral over the degree-$\mu$ covering $T$ is $\mu$.
This proves \eqref{eq:jacobian} without importing a Hahn flatness or
Nullstellensatz theorem.

For the second assertion, $f(0)=0$ makes each ordinary composition with $f$
a convergent germ, and expansion in $E$ is support-controlled. On the adapted
cycle,
\[
 G(F\circ\Phi)=G(t^{m_1\rho}p_1,\ldots,t^{m_n\rho}p_n).
\]
Its nonconstant target Taylor monomials have zero integral against
$\bigwedge dp_j/p_j$ on the torus; its constant monomial is $G(0)$. Finite
covering degree and strong summability give \eqref{eq:target-normalization}.
\end{proof}

The ordinary degree/residue identity is classical. The point here is that it
survives for completely unrestricted radius-free positive Hahn perturbations
and follows directly from an actual admissible parameterized cycle. Calling it
conservation of local algebra length would additionally require the relevant
finite-flatness identification; that identification is not smuggled into this
proof.

\subsection{A finite-sheet trace formula}
Let $\xi$ denote coordinates on $A$. For a function on the finite cover, its
fiberwise sum is a single-valued holomorphic function on $A$, coefficientwise.

\begin{corollary}[Protected trace representation]\label{cor:trace}
For $H\in\Agerm_n$, define
\[
 \Theta_H(\xi)=\sum_{q\in p^{-1}(\xi)}
 \frac{H(\Phi(q))}{J_F(\Phi(q))}.
\]
The finite sum descends to a Hahn-holomorphic function on $A$, and
\begin{equation}\label{eq:trace}
 \mathcal R_F(H)=\frac1{(2\pi i)^n}
 \HInt_{\T^n}\Theta_H(\xi)
 \frac{d\xi_1}{\xi_1}\wedge\cdots\wedge\frac{d\xi_n}{\xi_n}.
\end{equation}
\end{corollary}
\begin{proof}
Invertibility of $D_zF(\Phi)$ and the differential of \eqref{eq:adapted}
give
\[
 \Phi^*\left(\frac{H\dz}{\prod_jF_j}\right)
 =\frac{H(\Phi)}{J_F(\Phi)}\bigwedge_j\frac{dp_j}{p_j}.
\]
Integration over a finite cover is integration of the fiberwise sum. This is
valid for each ordinary coefficient and hence for the Hahn series. The
matrix inverse used here has a fixed finite monomial shift times a
positive-support expansion, so no summability issue is added.
\end{proof}

This is a trace over protected regular fibers, not an unsupported claim that
all individual zeros at target value $0$ have been separated or parameterized.

\subsection{Annihilation of the equation ideal}
\begin{proposition}\label{prop:annihilate}
For every $H\in\Agerm_n$ and each $j$,
$\mathcal R_F(F_jH)=0$. Thus $\mathcal R_F$ descends to a nonzero
$K$-linear functional on $\Agerm_n/(F_1,\ldots,F_n)$.
\end{proposition}
\begin{proof}
Classical residue annihilation gives $R_{f,k}(f_jh)=0$ when $k_j=0$, and
$R_{f,k}(f_jh)=R_{f,k-e_j}(h)$ when $k_j>0$. Insert these identities in
\eqref{eq:residue}, shift the $j$th summation index, and obtain
$\mathcal R_F(f_jH)=-\mathcal R_F(E_jH)$. All reindexings are allowed by
positive support. Nonzero follows from \eqref{eq:jacobian} in characteristic
zero. This proves descent, but not by itself nondegeneracy of a duality pairing.
\end{proof}

\section{A separately justified extension to formal coefficient germs}
\label{sec:formal}

The radius-free analytic theorem does not require a common radius, but it
still assumes that each coefficient germ is convergent. There is a genuine
extension beyond that hypothesis. The leading map $f$, however, remains an
ordinary convergent map with an isolated zero throughout this section.

\begin{lemma}[Finite-jet extension of ordinary residues]\label{lem:finitejet}
For fixed $f,k$, the functional $R_{f,k}$ depends on a finite Taylor jet of
its numerator. It has a unique extension to $\C[[z]]$ that is continuous
for the formal maximal-ideal filtration.
\end{lemma}
\begin{proof}
The ideal
\[
 I_k=(f_1^{k_1+1},\ldots,f_n^{k_n+1})\subset\C\{z\}
\]
is primary to the maximal ideal: its radical is the same as that of $(f)$.
Thus $\m^{a_k}\subset I_k$ for some finite integer $a_k$. Ordinary residue
annihilation implies that $R_{f,k}$ vanishes on $I_k$, and consequently on
$\m^{a_k}$. Evaluation on the Taylor polynomial of degree less than $a_k$
defines the extension. A different sufficiently high truncation has the same
value. Uniqueness follows because these polynomials are dense in the formal
maximal-ideal filtration. This is the classical finite-jet property of local
residue functionals, not a new duality theorem \cite{GH,TN}.
\end{proof}

For $E,H$ with formal coefficient germs, define $\mathcal R_F(H)$ by
\eqref{eq:residue}, using these extended $R_{f,k}$. Positive Hahn support
still makes that definition strongly summable.

\begin{theorem}[Formal-coefficient contour theorem]\label{thm:formal}
In \cref{thm:main,thm:adapted-main}, the assumptions
$E\in\Agerm_n^n$ and $H\in\Agerm_n$ may be replaced by
$E\in\FF_n^n$ and $H\in\FF_n$. The leading map $f$ is still convergent
and has an isolated zero. The same constructions and support certificates
apply, the pullback coefficients are ordinary holomorphic functions on the
fixed parameter domain, and the untransformed contour formula holds with the
finite-jet definition of the formal residue.
\end{theorem}
\begin{proof}
The construction of $p,X,T,Z$ uses only $f$, so it is unchanged. Every
substitution of a formal coefficient germ into $Z_0$, $\Phi$, or the
homotopy is justified by \cref{lem:substitution}. The implicit-lift proof
uses only those substitutions, derivatives of formal germs, and finite
coefficient operations, so its proof and support estimates also remain valid.
It remains to prove the residue equality; the ordinary constant-period lemma
cannot simply be quoted for a divergent formal numerator.

Fix an output Hahn exponent $\nu$. In the protecting-cycle integral, the
support argument \eqref{eq:exponent-accounting} admits only finitely many choices
of numerator exponent, perturbation exponents, and radial degrees. Each such
radial coefficient uses finite Taylor jets of the selected formal germs.
On the residue-series side, the support lemma admits only finitely many
coefficient choices and multi-indices $k$ at exponent $\nu$.
\Cref{lem:finitejet} shows that their residues also use finite Taylor jets.
Take the union of these two finite lists of jet requirements.

Replace each selected formal coefficient germ by a Taylor polynomial meeting
all its requirements. Leave the other exponents absent, or replace their
coefficients by zero. This replacement does not introduce new Hahn exponents
and does not alter the coefficient at $\nu$ on either side. The analytic
contour theorem, already proved, applies to these polynomial coefficient data
and identifies the two coefficients at $\nu$. Since $\nu$ was arbitrary,
the formal protecting-cycle identity follows.

Finally, the support-controlled homotopy proof is coefficientwise ordinary
Stokes on polynomial expressions in the real homotopy parameter; it never
uses convergence of an original coefficient germ. It identifies the adapted
period with the protecting period, proving the remaining identity.
\end{proof}

The finite-jet argument is a bridge between two specified rings, not an
identification of those rings. In particular, it does not assert that a
formal germ is an ordinary function on a complex neighborhood. It becomes a
Hahn family of holomorphic functions only after positive-support substitution.
Nor does the theorem construct the initial ramified covering when $f$ itself
is an arbitrary formal map. That is a different extension problem.

\section{A mixed six-sheeted example}
\label{sec:example}

\subsection{A singularity that is not already a coordinate product}
Work first over $\Gamma=\Q$ and set
\begin{equation}\label{eq:examplef}
 f_1=x^2+y^3,\qquad f_2=x^2+3y^3.
\end{equation}
The ideal $(f_1,f_2)$ equals $(x^2,y^3)$, so its colength is $\mu=6$.
The ordered denominator transformation has matrix
\[
 B=\begin{pmatrix}1&1\\1&3\end{pmatrix},\qquad \det B=2.
\]
It follows already at the ordinary origin that
$\Res_0(xy^2\,dx\wedge dy/(f_1f_2))=1/2$.
This factor will provide a useful check that the integrand has not been
silently replaced without its correct transformation factor.

For a target point $w=(w_1,w_2)$, the inverse equations are
\begin{equation}\label{eq:inverseexample}
 x^2=(3w_1-w_2)/2,\qquad y^3=(w_2-w_1)/2.
\end{equation}
The branch divisor lies in $(3w_1-w_2)(w_2-w_1)=0$. Choose target weights
$(1,2)$ and radial ramification $N=6$, so $(m_1,m_2)=(6,12)$. On a
polyannulus in variables $(a,b)$, put
\begin{equation}\label{eq:coverexample}
 p(a,b)=(a^6,b),
\end{equation}
a cover of degree $6$. The ordinary family may be written
\begin{align}
 Z_x(a,b,s)
 &=\sqrt{\frac32}\,s^3a^3
   \left(1-\frac{s^6b}{3a^6}\right)^{1/2},\label{eq:Zx}\\
 Z_y(a,b,s)
 &=-2^{-1/3}s^2a^2
   \left(1-\frac{s^6b}{a^6}\right)^{1/3}.\label{eq:Zy}
\end{align}
The binomial branches equal $1$ at $s=0$. They are single-valued and
holomorphic for small ordinary $s$ on a fixed collar of $|a|=|b|=1$.
Direct substitution gives
\[
 f_1(Z)=s^6a^6,\qquad f_2(Z)=s^{12}b.
\]
As $a$ runs through the six choices above a fixed $a^6$, its powers $a^3$
and $a^2$ give every pair of square-root and cube-root choices exactly once.
Thus this is the full six-sheeted cycle, not one arbitrarily chosen branch.

\subsection{A deformation and its exact adapted cycle}
Let
\begin{equation}\label{eq:exampleF}
 F_1=x^2+y^3-t,\qquad F_2=x^2+3y^3-5t.
\end{equation}
Here $\delta=1$ is a support lower bound. The Jacobian is
\begin{equation}\label{eq:exampleJ}
 J_F=\det\begin{pmatrix}2x&3y^2\\2x&9y^2\end{pmatrix}=12xy^2.
\end{equation}
On the family \eqref{eq:Zx}--\eqref{eq:Zy}, $x$ has radial order $3$
and $y$ has order $2$. The inverse Jacobian has pole order at most $4$.
Distinct sheets separate by radial order at most $3$. Thus one may take
\[
 b_{\mathrm{pole}}=4,\qquad c=3,\qquad L=13,\qquad \rho=\frac1{40}.
\]
The symbol $b_{\mathrm{pole}}$ here denotes the bound called $b$ in
\cref{lem:bounds}, not the second annular coordinate. The protecting inequality
is $17/40<1$. Write $s=t^{1/40}$, so $t=s^{40}$.

The adapted roots are characterized by
\begin{align}
 \Phi_x^2&=\frac{3s^6a^6-s^{12}b}{2}-s^{40},\label{eq:Phix}\\
 \Phi_y^3&=\frac{s^{12}b-s^6a^6}{2}+2s^{40},\label{eq:Phiy}
\end{align}
with the branches near $Z_x,Z_y$. Consequently
\[
 F_1(\Phi)=s^6a^6,\qquad F_2(\Phi)=s^{12}b
\]
exactly. These equations may be used directly as finite exact descriptions of
the two Hahn roots; the correct branches are selected by their leading terms.
Their first displacement terms are
\begin{align}
 \Phi_x-Z_x
 &=-\frac{s^{37}}{2\sqrt{3/2}\,a^3}+\smallO(s^{43}),\label{eq:dx}\\
 \Phi_y-Z_y
 &=\frac{2^{5/3}}{3a^4}s^{36}+\smallO(s^{42}).\label{eq:dy}
\end{align}
Here the $\smallO$ notation means a lower bound on the integer radial
support, not a limit in the surreal topology. The general estimate
$v(U)\geq\delta-b_{\mathrm{pole}}\rho=36/40$ is attained by the second
coordinate. The common ball $v(U)>13/40$ is much larger than the actual
correction; it is a uniqueness region, not a claim of a sharp bound.

\subsection{All monomial residues in closed form}
The displaced zeros of \eqref{eq:exampleF} satisfy
$x^2=-t$ and $y^3=2t$. They are six distinct algebraic points in an
algebraically closed Hahn workspace. The residue of a monomial is
\begin{equation}\label{eq:monomials}
 \boxed{\quad
 \mathcal R_F(x^a y^b)=
 \begin{cases}
 \dfrac12(-t)^p(2t)^q,& a=2p+1,\ b=3q+2,\ p,q\in\N,\\[2pt]
 0,&\text{otherwise},
 \end{cases}\quad}
\end{equation}
for nonnegative integers $a,b$.

To prove the formula independently, transform the ordered equations to
$(x^2+t,y^3-2t)$ and include the factor $1/2$. The one-variable residues
of $x^a/(x^2+t)$ and $y^b/(y^3-2t)$ are respectively $(-t)^p$ for
$a=2p+1$ and $(2t)^q$ for $b=3q+2$, and vanish otherwise. Equivalently,
use the quotient basis
\[
 1,\ y,\ y^2,\ x,\ xy,\ xy^2
\]
and reduce by $x^2=-t$, $y^3=2t$: the residue is one half of the
coefficient of $xy^2$. Either calculation agrees with the untransformed
cycle theorem, but neither substitutes for its arbitrary-support proof.
In particular,
\[
 \mathcal R_F(xy^2)=\tfrac12,\quad
 \mathcal R_F(x^3y^2)=-\tfrac t2,\quad
 \mathcal R_F(xy^5)=t,\quad
 \mathcal R_F(J_F)=6.
\]
The last identity checks the degree of the whole cycle and detects the loss
of a sheet or an incorrect orientation.

For the entire numerator $H=e^{x+y}$, the positive substitution is legitimate
and yields the exact Hahn series
\begin{equation}\label{eq:expresidue}
 \mathcal R_F(e^{x+y})=
 \frac12\sum_{p,q\geq0}
 \frac{(-1)^p2^q}{(2p+1)!(3q+2)!}\,t^{p+q}.
\end{equation}
This formula is a strong Hahn sum; its ordinary numerical convergence at a
complex value of $t$ is irrelevant to the theorem. The constant coefficient
is $1/4$, not $1/2$, because the $xy^2$ coefficient of the numerator is
$1/2$. Exact further coefficients are recorded by the accompanying check
program.

\subsection{Why the original coordinate torus is the wrong shortcut}
A tempting but incorrect replacement is to integrate the original form over
an arbitrary small product circle in $(x,y)$. On a torus where
$|y|^3\ll |x|^2$, both denominators of \eqref{eq:examplef} are dominated
by $x^2$; on a torus with the opposite inequality, both are dominated by
$y^3$. These cycles are not the equation-adapted Leray cycle.
For example, in the first region,
\[
 \frac{xy^2}{(x^2+y^3)(x^2+3y^3)}
 =x^{-3}y^2\sum_{r,s\geq0}(-1)^r(-3)^s
    \frac{y^{3(r+s)}}{x^{2(r+s)}}.
\]
There is no $x^{-1}y^{-1}$ term, so its coordinate-torus integral is $0$,
whereas the residue is $1/2$. The ramified cover constructed above fixes
this geometric problem without multiplying the original integrand by an
auxiliary determinant. A Jacobian from the actual parameterization remains,
as it must for every pullback of a differential form.

\section{Radius collapse and non-Archimedean support examples}
\label{sec:pathologies}

\subsection{Infinitely shrinking coefficient radii}
Keep \eqref{eq:examplef}, but replace the perturbation by
\begin{equation}\label{eq:radiuscollapse}
 E_1(x,y)=-\sum_{j\geq1}\frac{t^j}{1-jx},\qquad E_2=0.
\end{equation}
Its support is the positive integers, but the $j$th coefficient has a pole at
$x=1/j$. Therefore there is no common ordinary neighborhood of the origin
on which all coefficients are holomorphic. A proof based on a uniformly
convergent family on a fixed original source disk does not apply.

Nevertheless the protecting cycle \eqref{eq:Zx}--\eqref{eq:Zy}, with
$s=t^{1/40}$, still works. On that cycle,
\[
 \frac{1}{1-jZ_x}=\sum_{r\geq0}j^r Z_x^r
\]
is a strongly summable Hahn identity. Each exact output exponent sees only
finitely many $j,r$ and finitely many radial coefficient choices. The
holomorphic coefficient functions are all defined on the same annular
parameter domain because each is a finite polynomial expression in the
coefficients of $Z_x$. Their sizes may grow without bound as the exponent
varies; no coefficient norm bound was imposed or used.

The exact adapted lift also exists by \cref{prop:lift}, even though a finite
formula like \eqref{eq:Phix}--\eqref{eq:Phiy} is no longer available.
Its coefficients are obtained from finite rooted trees, all on the same
parameter domain, with the explicit support certificate \eqref{eq:Usupport}.

\subsection{A value group where radial powers are not cofinal}
Take $\Gamma=\Q^2$ with lexicographic order. For
$\rho=(0,1)$ and $\eta=(1,0)$ one has $r\rho<\eta$ for every positive
integer $r$. Hence no sequence of radial Taylor truncations indexed by
$r\in\N$ is automatically a valuation approximation at all scales. Yet
$\sum_{r\geq0}t^{r\rho}$ is a perfectly valid Hahn sum. Its support is
well ordered, and each exact exponent has a single radial degree.

One can combine this behavior with radius collapse and non-finitely-generated
support. For instance, use the well-ordered positive set
\[
 S=\{(0,1-1/j):j\geq2\}\cup\{(1,0)\}
\]
as the perturbation support and assign to the first set the coefficient germs
$(1-jx)^{-1}$. This support has order type $\omega+1$, has least element
$\delta=(0,1/2)$, and is not required to sit in a finitely generated
positive semigroup. Choosing $\rho=\delta/40$ supplies the same protecting
and lifting inequalities as before. The support proofs remain valid even
though infinitely many radial degrees lie below $(1,0)$.

\subsection{Why positivity cannot simply be removed}
The perturbation condition is structural. If arbitrary Hahn perturbations
were permitted, one could take $E=-f$, giving $F=0$ and no denominator at
all. More subtly, a support set unbounded below can destroy the finite-jet
property of substitution: the formal-looking series
$\sum_{r\geq0}t^{-r^2}z^r$ is not an element of $\FF_1$ with the required
well-ordered coefficient support. Such an expression must not enter the
argument merely because every individual coefficient exists.

The theorem does not assert that positivity is necessary for every individual
example. It gives a uniform sufficient condition under which the cycle,
substitution, inverse, and homotopy constructions all carry verifiable
support certificates.

\section{Construction, reproducibility, and proof obligations}
\label{sec:implementation}

\subsection{A finite-data construction guide}
When the initial germ is polynomial or effectively algebraic, the proof gives
an actionable route. Compute a finite-map representative or a local
elimination description, find a nonzero discriminant equation, and choose a
positive integer weight exposing one of its monomials. Compute the radial
monodromy or use an exponent dividing
$\operatorname{lcm}(1,\ldots,\mu)$ that kills it. Retain the finite
annular cover rather than forcing global single-valued branches on the target.
Expand its source-coordinate map in the ramified radial variable.

For a protecting cycle, only $m_*\rho<\delta$ is required. For an adapted
cycle, compute an inverse-Jacobian pole bound $b$ and a sheet-separation bound
$c$, choose $L$ and $\rho$ as in \eqref{eq:Lrho}, and solve
\eqref{eq:Bequation} by coefficient recursion. The support certificate
\eqref{eq:P} is part of the output. A requested individual Hahn coefficient
uses only finitely many tree terms; the theorem does not guarantee that an
arbitrary representation of its support makes those terms algorithmically
findable.

This distinction matters for computer algebra. Existence of a finite
coefficient dependency set does not by itself supply a decision procedure for
membership in an arbitrary well-ordered subset of a general ordered group.
Effective implementations should first use rational lattices, explicitly
presented support semigroups, algebraic coefficient germs, and exact branch
specifications. Arbitrary convergent germs need not be supplied by finite,
effective descriptions. No uniform complexity bound for all such inputs is
claimed.

\subsection{Executable finite checks}
The archive contains \texttt{code/checks.py}; the recorded run passes
632 exact checks. It uses exact rational
arithmetic and symbolic matrices, not floating-point sampling. Its tests
verify the six-dimensional quotient relations, the full monomial formula on
a finite grid, the Jacobian normalization, ideal annihilation on a finite
family, the target equations of both parameterizations, and finite binomial
jet identities. It also records coefficients of \eqref{eq:expresidue} and
the leading displacement constants. The output is stored in
\texttt{data/checks.json}; the script exits with an error if an asserted
identity fails.

The matrix check is deliberately different from the derivation of the
contour. Let $M_x,M_y$ be the multiplication matrices on the basis
$(1,y,y^2,x,xy,xy^2)$ over $\Q(t)$. The matrix
$12M_xM_y^2$ is invertible there, and
\[
 \operatorname{tr}\bigl(M_x^aM_y^b(12M_xM_y^2)^{-1}\bigr)
\]
computes the sum of the simple-zero residues of $x^a y^b$. Comparing it to
\eqref{eq:monomials} checks both the determinant factor and all six sheets.
It does not prove the arbitrary-support contour theorem. In particular, no
finite symbolic run can verify well-ordering for every input support,
parameterized removable singularities for every finite germ, or the general
rooted-tree construction.

\subsection{Dependency and audit table}
\begin{center}\small
\begin{longtable}{@{}L{0.25\textwidth}L{0.34\textwidth}L{0.33\textwidth}@{}}
\toprule
Step & Input or established result & What must be checked here\\
\midrule\endhead
Finite ordinary map & Isolated convergent zero; finite-map theorem
 & Proper representative, degree $\mu$, nonzero discriminant\\[0.4em]
Target weighting & Convergent nonzero discriminant
 & Unique least monomial and uniform annular avoidance\\[0.4em]
Radial ramification & Finite covering monodromy
 & Kill the radial permutation, not all angular monodromy\\[0.4em]
Extension across $s=0$ & Bounded source representative
 & Holomorphic parameter extension and $Z(q,0)=0$\\[0.4em]
Radius-free substitution & Positive well-ordered support
 & Finitely many terms at each exact exponent\\[0.4em]
Protecting cycle & $m_*\rho<\delta$
 & Positive denominator errors and certificate \eqref{eq:P0}\\[0.4em]
Residue equality & Classical fixed-germ Laurent period identity
 & Strong regrouping, with no common radius for all germs\\[0.4em]
Adapted cycle & Inverse-Jacobian bound $b$; choice $L>b$
 & Positive implicit coefficients, nullary-vertex support bound\\[0.4em]
Homotopy & $L>m_*$ and positive support
 & No denominator zero; coefficientwise closedness and Stokes\\[0.4em]
Formal extension & Residue finite jets
 & Separate finite-jet comparison, not an analytic assertion about a formal germ\\
\bottomrule
\end{longtable}
\end{center}

\section{Position in the literature and remaining questions}
\label{sec:scope}

\subsection{What is classical and what is proposed as new}
Ordinary Leray cycles and the Grothendieck residue are classical
\cite{GH,TN}. So are finite holomorphic maps, analytic discriminants, and
Weierstrass preparation \cite{GR}. The passage from a finite unramified
cover to ramified single-valued analytic data is an established ingredient
of Abhyankar--Jung theory; Parusi\'nski and Rond explicitly use a finite
covering, power substitution, and removable singularities in their analytic
proof \cite{PR}. The positive support lemma is classical
\cite{Neumann,Higman}. None of these ingredients is claimed here as an
independent discovery.

The proposed contribution is their combination with arbitrary Hahn support
and radius-free coefficient germs to obtain an original-form period on a
single fixed finite-cover cycle, uniformly for all perturbations above a
specified positive valuation. The support bounds certify an exact adapted
lift and an admissible homotopy in that same category. The theorem includes
higher-rank value groups and formal coefficient perturbations by a separate
finite-jet argument. These qualifications are not supplied merely by naming
the ordinary residue theorem or the analytic Abhyankar--Jung theorem.

M\"uller--Strohmaier's Hahn-holomorphic and Hahn-meromorphic functions provide
another important use of Hahn series in complex analysis \cite{MS}. Their
analytic summation setting should not be identified with the present
coefficientwise ring of arbitrary radius-free germs. Likewise, the
repository's determinant-transformed formula remains valid and useful; the
new formula answers its explicitly different original-integrand question.
The new cycle need not be a coordinate torus, and its finite angular cover
is essential in general.

The literature check was targeted, performed on September 21, 2026, and
included the repository's contour problem, ordinary residue mappings,
finite-map ramification, Hahn-meromorphic functions, and searches combining
Hahn series with Leray cycles and Grothendieck residues. No matching theorem
with the exact arbitrary-support and radius-free hypotheses was located.
Failure to locate a theorem does not certify historical priority. The claim
that can be evaluated directly is the stated theorem and its proof, not a
universal assertion about all published literature.

\subsection{Exact boundaries of the result}
All sums, supports, covers, and parameter manifolds are set-sized. Every
coefficientwise integral is an ordinary integral on a compact real manifold.
The Hahn parameterization is not asserted to be continuous for the fine
surreal topology, and the homotopy is not a fine-topological homotopy. No
surreal-valued measure or unrestricted Riemann-sum construction is invoked.

The input must have an isolated convergent ordinary leading map and a
positive-support perturbation. The result does not handle a completely
arbitrary Hahn tuple with no such leading map, or a non-isolated leading zero.
It does not produce a canonical contour independent of choices; it proves
that the resulting period is independent of those choices. It does not
assert that the parameter space is connected, or that every possible cycle
is homologous to the chosen one.

The adapted parameterization gives $\mu$ distinct regular points over each
protected target point. It does not assert that every zero of $F$ is simple.
The residue is well defined even when the deformed zero algebra is not
reduced. Nonzero descent of the residue functional is established, but a full
perfect-pairing or flatness theorem is not inferred from it without proof.

\subsection{Further problems now made more precise}
One natural next question is whether the convergent leading map can itself be
replaced by an arbitrary formal isolated map while retaining a fixed ordinary
compact parameter manifold. The present proof depends on a bounded analytic
finite covering before Hahn substitution; formal algebra alone does not
supply that object.

A second question concerns effective scale selection. The estimates in
\eqref{eq:Lrho} are transparent and sufficient, but can be far from sharp.
An intrinsic invariant of the finite map controlling the best protecting and
lifting scales could improve both the mathematics and exact symbolic
implementations. The attained bound $36/40$ in the six-sheeted example shows
that the inverse-Jacobian pole order carries meaningful information, even
when the general uniqueness radius is conservative.

A third question asks for a homology comparison among different protecting
weights, covers, and microscopic scales. Equality of their residue periods
is proved here. A chain-level comparison for a precisely specified larger
Hahn category would be stronger and is not a consequence of period equality
alone.

\section{Conclusion}
The original-integrand cycle problem has a concrete solution in the
finite-cover, coefficientwise Hahn chain category. A generic target monomial
direction avoids the ordinary discriminant; finite radial ramification gives
a common parameter space; a sufficiently small positive exponent gives a
protecting infinitesimal radius; and the positive support lemma turns every
radius-free substitution into finite coefficient algebra. The resulting
integral equals the independently defined perturbation residue.

The exact adapted lift goes further: its target equations hold identically,
its sheets remain separated, and a support-controlled pole-free homotopy
connects it to the universal protecting cycle. The proof therefore supplies
not merely an existence assertion, but an explicit support certificate,
normalization identities, a regular-fiber trace description, and a formal-coefficient extension. The surviving limitations concern different categories
and genuinely additional hypotheses, rather than an omitted common-radius
or summability assumption.

\begin{thebibliography}{99}
\small\raggedright
\bibitem{Repo}
VladimirReshetnikov/Surreal contributors,
\emph{Surreal repository: the surreal and surcomplex reports},
repository catalogue and source files, inspected September 21, 2026,
commit \texttt{39f2be6667ade51bca2b45daa47e289d69c09764}.
\url{https://github.com/VladimirReshetnikov/Surreal}.
The catalogue labels the research reports AI-assisted, unrefereed drafts.

\bibitem{RepoContours}
\emph{Jordan Separation, Cauchy Theory, and Stokes Calculus on the Surcomplex
Plane}, merged research report in \cite{Repo}, September 21, 2026.
See the scope discussion after the general determinant-transformed residue
representation and the subsection ``Research problems sharpened by these
results.'' Source path:
\path{docs/surcomplex/contours-and-stokes/article.tex}.
\url{https://github.com/VladimirReshetnikov/Surreal/tree/39f2be6667ade51bca2b45daa47e289d69c09764/docs/surcomplex/contours-and-stokes}.

\bibitem{RepoAnalysis}
\emph{Surcomplex analysis report}, in \cite{Repo}, source directory
\path{docs/surcomplex/analysis}, September 21, 2026.
The radius-free germ and finite-deformation companion reports are in
\path{docs/surcomplex/analytic-geometry} and
\path{docs/surcomplex/finite-deformations}.

\bibitem{GH}
P.~Griffiths and J.~Harris,
\emph{Principles of Algebraic Geometry}, Wiley, 1978; Wiley Classics Library
edition, 1994. Chapter~5, ``Residues,'' pp.~647--732.
\url{https://doi.org/10.1002/9781118032527}.

\bibitem{TN}
S.~Tajima and K.~Nabeshima,
\emph{An effective method for computing Grothendieck point residue mappings},
arXiv:2011.09092, 2020.
\url{https://arxiv.org/abs/2011.09092}.
See the definition of the local residue, transformation law, and finite local
algebra discussion.

\bibitem{GR}
H.~Grauert and R.~Remmert,
\emph{Coherent Analytic Sheaves}, Grundlehren der mathematischen Wissenschaften
265, Springer, 1984. Chapters~II--III: local Weierstrass theory and finite
holomorphic maps.
\url{https://doi.org/10.1007/978-3-642-69582-7}.

\bibitem{PR}
A.~Parusi\'nski and G.~Rond,
The Abhyankar--Jung theorem,
\emph{Journal of Algebra} \textbf{365} (2012), 29--41.
\url{https://doi.org/10.1016/j.jalgebra.2012.05.003};
\url{https://arxiv.org/abs/1103.2559}.
See in particular the analytic covering argument and its ramification and
removable-singularity steps; the ordinary mechanism is classical.

\bibitem{Neumann}
B.~H.~Neumann,
On ordered division rings,
\emph{Transactions of the American Mathematical Society} \textbf{66} (1949),
202--252.
\url{https://doi.org/10.1090/S0002-9947-1949-0032593-5}.

\bibitem{Higman}
G.~Higman,
Ordering by divisibility in abstract algebras,
\emph{Proceedings of the London Mathematical Society}, series~3,
\textbf{2} (1952), 326--336.
\url{https://doi.org/10.1112/plms/s3-2.1.326}.

\bibitem{MS}
J.~M\"uller and A.~Strohmaier,
The theory of Hahn-meromorphic functions, a holomorphic Fredholm theorem,
and its applications,
\emph{Analysis \& PDE} \textbf{7} (2014), no.~3, 745--770.
\url{https://doi.org/10.2140/apde.2014.7.745};
\url{https://arxiv.org/abs/1205.0236}.
\end{thebibliography}
\end{document}
